% Chapitre « Hyperparamètres » (axe 1). Sources : v2/docs/lab/fiches/param-*.md (29 fiches, sections
% « Relecture (2026-11-15) »), v2/docs/lab/reports/2026-11-14-axe1-hyperparametres.md,
% v2/docs/lab/reports/2026-11-15-axe1-relecture-sources.md et son dossier (JSON des lecteurs et vérificateurs),
% v2/docs/lab/reports/2026-11-14-audit-sources.md, v2/papers/axe1-hyperparametres/main.tex. État relu le 2026-11-15.
\providecommand{\donfic}[1]{\nolinkurl{#1}}
\section{Hyperparamètres}\label{par:sec}\label{sec:parametres-ref}
\begingroup\setlength{\emergencystretch}{3em}% local à ce chapitre

Un hyperparamètre est une valeur que le jumeau ne mesure pas sur ses propres séries : une règle de régulation,
une part de filière, un taux d'abandon en formation, un âge d'entrée, un flux de diplômés étrangers. Chacun a
une fiche \donfic{v2/docs/lab/fiches/param-*.md} qui donne sa valeur, sa source primaire, son millésime et son
statut. Abréviations : \donfic{R/axe1} pour \donfic{2026-11-14-axe1-hyperparametres.md}, \donfic{R/relecture}
pour \donfic{2026-11-15-axe1-relecture-sources.md}, \texttt{Audit} pour \donfic{2026-11-14-audit-sources.md},
\texttt{A1} pour l'article \donfic{v2/papers/axe1-hyperparametres/main.tex}, \texttt{D} pour
\donfic{v2/docs/lab/decisions.md}.

\paragraph{Convention.} La convention est celle de \texttt{A1}. \mesure{}~: calculé par le projet sur une
donnée ouverte, ou lu dans la source primaire citée à côté du nombre (statut « sourcé », So, des fiches).
\derive{}~: rapport, différence ou arrondi calculé par nous. \suppose{}~: valeur par défaut sans mesure.
Statuts des fiches : So sourcé, Me mesuré, Su supposé, NM non mesurable. Les années, dates, numéros
d'articles et de pages ne sont pas étiquetés.

\subsection{Catalogue}\label{par:catalogue}

Le tableau~\ref{par:tab-catalogue} reprend les 29\mesure{} fiches. La dernière colonne donne le verdict de la
relecture en texte intégral du 2026-11-15 (section~\ref{par:relecture}) : \emph{confirmé}, \emph{corrigé de
$\to$ à}, ou \emph{introuvable}. Une fiche « confirmée » peut avoir reçu une précision de libellé sans
changement de valeur.

{\scriptsize
\setlength{\tabcolsep}{2pt}
\begin{longtable}{@{}p{2.25cm}p{4.0cm}p{1.55cm}p{2.75cm}p{0.6cm}p{3.75cm}@{}}
\caption{Catalogue des hyperparamètres et résultat de la relecture du 15/11 (fiches param-$\ast$,
\donfic{R/relecture}, tableau « Résultat »).}\label{par:tab-catalogue}\\
\toprule
Fiche & Valeur & Unité & Source & Statut & Relecture du 15/11 \\
\midrule\endfirsthead
\toprule
Fiche & Valeur & Unité & Source & Statut & Relecture du 15/11 \\
\midrule\endhead
\bottomrule\endfoot
\multicolumn{6}{@{}l}{\emph{Régulation de l'installation}}\\
ide-regul-regle & zone « sur-dotée » (terme de l'avenant), dite ici zone remplacement seul : 1\mesure{} entrée pour 1\mesure{} cessation ; en périphérie, au moins
2/3\mesure{} de l'activité dans la zone ; 27,4~\%\derive{} de la population, 13,7~\%\derive{} des communes &
règle par commune & avenant 6 IDE, art. 3.4 \citep{d-avenant6ide} ; zonage au 20/10/2026
\citep{d-zonagesdatagouv} & So & confirmé \\
ide-regul-effet-densite & $-25{,}57$\mesure{} IDEL pour \num{100000} hab. ($-18{,}8$\mesure{} chez les moins de
40~ans), $p<0{,}001$\mesure & différence de différences & IRDES QES 270, p.~5, 2006--2016 \citep{d-irdes270} &
So & confirmé ; appariement p.~3 ; repères 2012 sourcés par la presse \\
regul-theta-fuite & $\theta = \num{0.80}$ $[\num{0.74} ; \num{0.87}]$\suppose & sans dimension & rapport interne
axe~3, N3 (tendances parallèles : $p = \num{0.043}$\mesure) & Su & confirmé (rapport interne) \\
sf-regul-regle & 1 pour 1\mesure{} en zone remplacement seul ; 12,2~\%\derive{} de la population, 4,2~\%\derive{} des
communes & règle & avenant 4 SF, art. 3.3 \citep{d-avenant4sf} & So & corrigé (mécanisme omis) : les entrées
dérogatoires (art. 3.3.3 a-b) ne rendent pas de place \\
sf-regul-effet & direction seulement (disparités en baisse, 2016--2019) & --- & IGAS 2022-027R, p.~67
\citep{d-igas2023} ; Cour des comptes 2024 \citep{d-ccomptes2024} & NM & corrigé (attribution) : le bilan GDR
porte sur les aides, pas sur la régulation \\
mk-regul-regle & 1 pour 1\mesure{} en zone non prioritaire ; cédant $\ge \num{1200}$\mesure{} actes ; successeur
sous 2\mesure{} ans ; 26,5~\%\derive{} de la population, 8,1~\%\derive{} des communes & règle & avenant 7 MK,
art. 1.2.1 \citep{d-avenant7mk} ; CE n\textsuperscript{o}~488665 \citep{d-ce488665} & So & confirmé ;
« avenant 5 » non sourcé comme texte de régulation \\
mk-obligation-exercice & 2\mesure{} ans (au moins \num{6000}\mesure{} actes) en zone « sous-dotée » (terme de l'avenant), ou 2\mesure{}
ans en établissement ; premiers diplômés concernés vers 2027--2028\derive & années & avenant 7 MK, art. 4.1.1
\citep{d-avenant7mk} & So & confirmé \\
mk-regul-effet & non chiffré & --- & IRDES DRIMK (juillet 2026, p.~4) ; Cour des comptes 2024
\citep{d-ccomptes2024} & NM & corrigé : travail IRDES achevé, sens établi, ampleur non publiée ; « avenant 5
(2018) » $\to$ « accord conventionnel de 2018 » \\
cd-regul-regle & 1 pour 1\mesure{} en zone non prioritaire ; cédant $\ge 2$\mesure{} jours par semaine ;
successeur sous 1\mesure{} an ; 4,8~\%\mesure{} de la population ; spécialistes dérogatoires (art. 35.5) & règle
& convention CD 2023, art. 33.1, 35 \citep{d-conventioncd2023} ; ARS PACA \citep{d-arspaca2025cd} & So &
confirmé ; art. 35.5 relibellé « dérogation » dans le code \\
cd-regul-effet & non chiffré ; borne haute d'environ 5~\%\derive{} du flux & --- & mémoire EHESP, \S~3.3.2,
p.~53 \citep{d-bahroun2025} & NM & confirmé \\
regul-report-entrants ($\psi$) & qualitatif : report en bordure ; « contournement par le remplacement » &
parts par destination & IRDES QES 270, p.~6 \citep{d-irdes270} ; Cour 2024, p.~95 \citep{d-ccomptes2024} & NM
& introuvable : « remplacement » et « deux tiers (EHESP) » absents des deux sources ; report en bordure
confirmé \\
\midrule
\multicolumn{6}{@{}l}{\emph{Filières en pharmacie}}\\
ph-filiere-primo-inscrits ($\varphi$) & 2025 : D 71,0 ; H 9,5 ; B 8,6 ; G 5,7 ; E 3,5 ; A 1,2 ; C
0,4~\%\mesure{} de \num{3044}\mesure{} primo-inscrits & part par section & CNOP, Panorama 2025, p.~2
\citep{d-cnop2025} & So & confirmé (étiquettes vérifiées sur le rendu image) \\
ph-part-privee-biologie ($s_G$) & 55,8~\%\derive{} (2025) ; 55,3~\%\derive{} (2024) & part de la section G &
CNOP 2025 \citep{d-cnop2025} & So & confirmé ; pente $-0{,}5$\derive{} point par an plutôt $-0{,}4$\derive \\
ph-part-privee-etablissements ($s_H$) & 28,6~\%\derive{} des inscriptions H+E (\num{2548} / \num{8924}\mesure) ;
30,0~\%\derive{} des fonctions en PUI & part des inscriptions & CNOP 2016, 2024, 2025
\citep{d-cnop2016,d-cnop2024,d-cnop2025} & So & corrigé : « \num{8106} personnes » $\to$ \num{8106}\mesure{}
inscriptions, \num{7817}\mesure{} personnes ; ligne 2015 en double \\
ph-internat-postes & \num{746}\mesure{} postes (biologie \num{266}, pharmacie hospitalière \num{480}\mesure),
rentrée 2025 ; \num{536}\mesure{} internes inscrits en 1\textsuperscript{re} année ; environ 19~\%\derive{} d'une
promotion & postes par an & arrêté du 6/12/2024 \citep{d-arrete2024pharma} ; CNOP 2024 \citep{d-cnop2024} & So
& confirmé \\
ph-taux-inscription-diplomes ($\varepsilon_{\mathrm{PH}}$) & 32,8~\%\mesure{} de diplômés non inscrits à 3 ans
(2015) & part & CNOP 2016, p.~7 \citep{d-cnop2016} & So & confirmé \\
ph-sorties-ordre & \num{3902}\mesure{} radiations en 2024 ; 18,4~\%\mesure{} de 60~ans et plus & personnes par
an & CNOP, synthèse 2024 \citep{d-cnop2024} & So & corrigé : ajout du solde primo-inscrits $-$ radiations $=
\num{-975}$\derive \\
\midrule
\multicolumn{6}{@{}l}{\emph{Odontologie}}\\
cd-specialisation-des ($\sigma_{\mathrm{CD}}$) & 9,1~\%\derive{} d'une promotion (2024 : 121 /
$\approx\num{1332}$) ; ODF 52, CO 25, MBD 44\mesure{} ; 140\mesure{} postes en 2026-27 & postes / places en
2\textsuperscript{e} année à $t-4$\suppose & arrêtés 2023, 2024, 2026
\citep{d-arrete2023odonto,d-arrete2024odonto,d-arrete2026odonto} ; CNG \citep{d-cng2018} & So & corrigé :
2023 = 120 / \num{1332} = 9,1 $\to$ 9,0~\%\derive{} ; 2024 confirmé \\
cd-stock-specialistes & 2022 : ODF \num{2465}\mesure{} (5,4~\%\derive), CO 222\mesure, MBD 100\mesure,
omnipraticiens \num{43242}\mesure & praticiens & ONCD 2019--2022 \citep{d-oncd2022} & So & corrigé : ODF
5,3 $\to$ 5,4~\%\derive{} (\num{2465} / \num{46029}) \\
\midrule
\multicolumn{6}{@{}l}{\emph{Formation infirmière}}\\
ide-rendement-diplomes ($\rho$) & \num{0.72}\mesure{} (\num{26020} / \num{36218}, entrée 2022) ; \num{0.81} à
\num{0.85}\mesure{} sur 2009--2018 & diplômés$(t+3)$ / inscrits$(t)$ & DREES 491 \citep{d-drees491} & Me &
confirmé (\num{0.72}) ; « 76~\% » = diplômés 2022 / inscrits 2019 ; « \num{0.77} » introuvable \\
ide-abandon-promotion & 25,0~\%\mesure{} de la promotion 2025 ; 11,9~\%\mesure{} définitives (promotion 2024) &
part d'une promotion & DREES 491 ; ER 1266 \citep{d-drees491,d-er1266} & So & confirmé ; « sans report »
jusqu'en 2021 \\
ide-abandon-1a & 12,0~\%\mesure{} (2025) ; 3,3 (2011) $\to$ 10,0 (2021) $\to$ 13,1 à 13,5~\%\mesure{}
(2022--2024) & part des inscrits de 1\textsuperscript{re} année & DREES 491 \citep{d-drees491} & So & confirmé ;
« quadruplé » $\to$ « triplé entre 2011 et 2021 » \\
\midrule
\multicolumn{6}{@{}l}{\emph{Sorties et entrée en exercice}}\\
cessation-par-age $x_p(a)$ & non publiée ; mesurable sur le RPPS & probabilité annuelle & DD76, tab.~11
\citep{d-dd76} ; ER 1062, 1075 \citep{d-er1062,d-er1075} & NM & corrigé (libellé) : pas de ventilation par
sexe pour les SF ; ancrages agrégés ajoutés \\
mortalite & incluse dans la cessation & --- & DD76 ; ER 1062, 1075, 1319 \citep{d-dd76,d-er1319} & NM &
confirmé \\
age-entree-exercice & diplômés en France : SF 24, CD 26, PH 27 à 29~ans\mesure{} ; à l'étranger : SF 28, CD
28~ans\mesure & âge moyen & DD76, tab.~12-13 \citep{d-dd76} & So & confirmé, complété : MK diplômés à l'étranger
26~ans\mesure{} \citep{d-er1075} \\
cumul-emploi-retraite & SF 0,8 ; CD 4,7 ; PH 1,8 ; MK et IDE 1,5~\%\derive & part des cotisants & COR 2016,
doc.~11 \citep{d-cor2016} & So & confirmé ; répartition MK / IDE de la CARPIMKO\suppose \\
mk-age-retraite & 66,8 à 67,3~ans\mesure{} (moyenne) ; 657 à 811\mesure{} départs par an & années & CNOMK 2024,
tab.~5, 7 \citep{d-cnomk2024} & So & confirmé ; libellé « cessation d'activité » précisé \\
\midrule
\multicolumn{6}{@{}l}{\emph{Ajoutés après l'audit du 14/11}}\\
part-privee-salaries-sae & 2025 : SF 19,8 ; PH 29,4 ; MK 44,3 ; CD 2,2~\%\mesure & part des salariés au 31/12 &
SAE, jeu 708, Q20, Q23 \citep{d-sae708} & Me & confirmé (codes ; parts non recalculées) \\
diplomes-etrangers-flux & CD \num{1300}\derive, MK \num{2200}\derive, IDE 400\mesure, PH 240\derive, SF
80\derive{} par an & personnes par an & DD76, ER 1026, 1075, 1319, CNOMK, CNOP
\citep{d-dd76,d-er1075,d-er1319,d-cnomk2024} & So & confirmé (valeurs) ; étiquettes corrigées : CD \num{1300}
sourcé $\to$ dérivé \\
\end{longtable}
}

\subsection{Relecture intégrale du 15 novembre}\label{par:relecture}

\paragraph{Protocole (\donfic{R/relecture}, \S~Comment).} 23\mesure{} lecteurs (modèles Haiku), un par source,
ont relu en texte intégral 355\mesure{} affirmations des fiches : 247\mesure{} confirmées, 35\mesure{} écarts,
67\mesure{} introuvables, 6\mesure{} inaccessibles. Chaque écart et chaque introuvable est passé par un
vérificateur (modèle Sonnet, 10\mesure{} agents) chargé de le réfuter : 6\mesure{} écarts réfutés,
27\mesure{} introuvables retrouvés, souvent dans une autre source de la même fiche. Les consignes et les JSON
sont dans \donfic{v2/docs/lab/reports/2026-11-15-axe1-relecture/}.

\paragraph{Verdict par paramètre.} Sur 29\mesure{} paramètres : \textbf{20\mesure{} confirmés, 8\mesure{}
corrigés, 1\mesure{} introuvable}. Nous avons recompté ces nombres sur le tableau « Résultat » de
\donfic{R/relecture} : les 8 corrigés sont cd-specialisation-des, cd-stock-specialistes, cessation-par-age,
mk-regul-effet, ph-part-privee-etablissements, ph-sorties-ordre, sf-regul-effet et sf-regul-regle ;
l'introuvable est regul-report-entrants. Aucune source n'était entièrement inaccessible ; trois l'étaient en
partie (base RegProf, communiqué DREES du 07/11/2023, Légifrance lu via des reproductions du JO).

\paragraph{Effet sur le code.} Aucune valeur numérique du code ne change. Les valeurs FLUX (\num{1300},
\num{2200}, 400, 240, 80), PART\_CUMUL, ABANDON (IDE : rendement $\rho = \num{0.718}$\mesure{}, soit $1-\rho^{1/3} \approx 10{,}5$~\%\derive{} d'abandon par an, \donfic{v2/jumeau/modules/formation.py} l.~22 ; SF \num{0.0331} et $L = 4$) et
$\theta = \num{0.80}$\suppose{} sont confirmées ou défendables. Seuls des libellés changent, dans
\donfic{regulation.py}, \donfic{etrangers.py} et \donfic{sorties.py}. \donfic{jumeau.check} passe
18\mesure{} critères sur 18, avec des sorties identiques à la référence.

\paragraph{Écart mesuré, non corrigé : l'âge d'entrée des diplômés à l'étranger.} Le code fixe un âge unique de
28\suppose{} ans. Les sources donnent 26~ans\mesure{} pour les MK \citep{d-er1075}, 31~ans\mesure{} pour les PH
\citep{d-dd76}, et pour les CD 32~ans\mesure{} (ER 1026) ou 28~ans\mesure{} (DD76) : elles se contredisent.
Dans le scénario « étrangers $\times 2$ », l'écart sur $N(2050)$ vaut 0\mesure{} pour les MK (même tranche
d'âge), $-191$\mesure{} ($-0{,}2$~\%\derive) pour les PH et $\num{-1369}$\mesure{} ($-1{,}3$~\%\derive) pour les
CD (\donfic{impact_age.py}). Le tendanciel ne bouge pas. C'est ouvert dans \texttt{D}.

\paragraph{Réserves.} Les lecteurs ont souvent lu par mots-clés ; environ la moitié\mesure{} des sources n'a eu
qu'un accès partiel. Les 247 affirmations « confirmées » n'ont pas eu de second passage : un faux confirmé n'a
pas été cherché. Deux sources se contredisent elles-mêmes : ONCD 2022 (\num{1333}\mesure{} diplômés hors France
p.~27, \num{1321}\mesure{} d'après le tableau p.~28) ; dénominateur d'odontologie (\num{1332}\mesure{} selon
SNESUP et Weka, \num{1322}\mesure{} selon L'Étudiant).

\subsection{Contradictions entre l'audit du 14 et la relecture du 15}\label{par:contradictions}

L'audit du 14/11 a relu 469\mesure{} citations du projet (\texttt{Audit}) ; la relecture du 15/11 a relu les
sources des fiches param-$\ast$. Elles se contredisent sur deux points. Dans les deux cas, la fiche suit la
relecture.

\paragraph{IGAS, note 133 (MK).} L'audit (registre A1, fiche \emph{mk-regul-effet}) écrit que l'IGAS ne
date pas la hausse des disparités d'accès aux MK (« note 133, sans période ; le \S~200 couvre 2016--2021 ») et
la classe CORRIGÉE (période). La relecture (fiche \emph{mk-regul-effet}, section « Relecture ») écrit que la
période est bien 2016--2019 et que « sans période précise » est faux \citep{d-igas2023}. Le texte de l'IGAS n'a
pas été rouvert une troisième fois.

\paragraph{Avenant 5 (2018) des MK.} L'audit (registre A1) ajoute « la régulation de l'avenant 5 (2018) en zones
sur-dotées », en s'appuyant sur la note 152 de la Cour des comptes (p.~95) et sur la question écrite AN 14922,
qui parle d'un dispositif « issu des avenants n° 5 et n° 7 ». La relecture lit la note 152 : la Cour écrit
« accord conventionnel de 2018 », sans numéro \citep{d-ccomptes2024}. L'avenant 7 ne cite l'avenant 5 qu'à
propos d'options conventionnelles \citep{d-avenant7mk}. Les deux vérificateurs concluent que le numéro n'est
pas sourcé ; l'un le marque dérivé, l'autre supposé (\donfic{verificateurs/v07}, \texttt{v08}). La question
écrite AN 14922 n'est pas dans le périmètre de la relecture : le point n'est donc pas clos. \texttt{A1} (\S~Régulation)
écrit encore « l'avenant~5 de 2018 » ; le tableau~3 de l'article de l'axe~4 et \donfic{v2/data/events.csv}
gardent une régulation MK au 01/07/2018, que l'audit (registre C2) jugeait déjà non vérifiable.

\subsection{Part des paramètres réellement sourcés}\label{par:sources}

Le décompte des statuts a changé entre deux documents : 27\mesure{} fiches dans \donfic{R/axe1} (19 So, 1 Me,
1 Su, 6 NM\mesure) ; 29\mesure{} dans \texttt{A1} et \donfic{R/relecture} (20 So, 2 Me, 1 Su, 6 NM\mesure),
après l'ajout de \texttt{part-privee-salaries-sae} (Me) et \texttt{diplomes-etrangers-flux} (So).

Le statut « sourcé » de 20\mesure{} fiches sur 29, soit 69~\%\derive, surestime ce qui est lu tel quel dans
une source :
\begin{itemize}
  \item 5\mesure{} fiches So ont été corrigées le 15/11 (sf-regul-regle, ph-part-privee-etablissements,
  ph-sorties-ordre, cd-specialisation-des, cd-stock-specialistes) : 15\derive{} fiches So sur 29, soit
  52~\%\derive, sont confirmées sans correction ;
  \item parmi elles, \texttt{diplomes-etrangers-flux} porte des valeurs presque toutes dérivées (CD, MK, PH,
  SF) ; seule la valeur IDE (400) est lue dans une source \citep{d-er1319} ;
  \item plusieurs fiches So portent une règle sourcée et une part de population dérivée (zonage) ; c'est la
  règle qui est sourcée, pas la part.
\end{itemize}
Au total, 14\derive{} fiches sur 29, soit 48~\%\derive, ont une valeur principale lue dans une source primaire
et confirmée sans correction. 2\mesure{} sont mesurées par le projet. Les 13\derive{} autres sont non
mesurables (6), supposée (1), corrigées (5) ou dérivées pour l'essentiel (1). Les ordres de grandeur du jumeau
reposent donc, pour moitié, sur des valeurs qui ne sont pas des lectures directes.
\endgroup
